\section[1-移位 Cartier 对偶]{1-移位 Cartier 对偶（3 月 23 日）}
\srcpage{41}

刚放完春假，先回顾上次见面时在做的事：把 Tannakian 重构解释成「取簇上的函数」与
「取交换环的谱」之间那个伴随的范畴化类似物。从这个视角，对仿射代数群 $G$，
我们把 $\Rep G$ 恢复为叠 $\mathrm{pt}/G$ 上拟凝聚层的范畴。

我们还进一步用谱分解（在 Fourier–Mukai 对偶、即 1-移位 Cartier 对偶的语境中）解释了它：
若 $G$ 是阿贝尔群，$\mu:G\times G\to G$ 是其乘法映射，则沿乘法的前推在 $\QCoh(G)$ 上
定义了卷积对称幺半积。例如取 $\QCoh(\Z)$（或 $\Z$-分次向量空间），其谱是 $B\G_m$，
因为这个范畴作为对称幺半范畴等价于 $\G_m$-表示的范畴。「移位」是因为我们从 $\Z$ 出发，
得到的却是 $B\G_m$ 而非 $\G_m$。

更一般地，设 $G$ 是有限阿贝尔群概形，令 $G^\vee:=\Spec\bigl(\C[G],*\bigr)$。
于是 $G$ 上带卷积的函数，等同于 $G^\vee$ 上带逐点乘法的函数。再高一格，
$\Rep G\simeq\QCoh(G^\vee)$。移位再次出现，因为 $\Rep G\simeq\QCoh(BG)$：
升到范畴层面时，总得把其中一边下移。

退一步问：$G^\vee[1]$，即 Fourier–Mukai 对偶，究竟是什么？它的点是什么？
我们是通过函子观点定义它的；对域 $k$，$G^\vee[1]$ 的 $k$-点是对称幺半函子
\[
\bigl(\mathrm{Shv}(G),*\bigr)\simeq\bigl(\QCoh(G),*\bigr)\longrightarrow\Vect_k .
\]
可以把它们想成「$(\QCoh(G),*)$ 的范畴化一维表示」——未范畴化的表示是映射
$G\to\mathrm{GL}_1(k)$，等价于代数同态 $\C[G]\to\End(k)$；
而这里我们有到 $\End(\Vect_k)=\Vect_k$ 的对称幺半函子。

因此 $G$ 应当映到 $\Aut(\Vect_k)$，即由线（在 $\otimes$ 下可逆的向量空间，
等价于一维向量空间）带张量积构成的对称幺半范畴。这个对象等价于 $B\G_m$，
是一个「群」，更确切地说是叠中的群对象。也就是说，它的函子观点取值不在阿贝尔群里
（如群概形那样），也不仅在群胚里（如通常的叠那样），而在 Picard 群胚里：
即所有对象都在 $\otimes$ 下可逆的对称幺半群胚。\yizhu{Picard 群胚的同构类构成一个阿贝尔群，
其自同构群给出稳定化子；它是「群」这个概念在范畴层面的推广。}{B（nLab "Picard groupoid"）}

再展开一点：映射 $X\to B\G_m$ 就是 $X$ 上线丛带张量积构成的对称幺半范畴，
而这些线丛在张量积下全都可逆。为什么是 $B\G_m$？思路是：所有可逆 $k$-向量空间都同构，
但它们的同构构成一个 $\G_m$，所以这个叠等价于 $\mathrm{pt}/\G_m$，
而这正是我们最初对 $B\G_m$ 的定义。

要点之一是：$\Spec\bigl(\mathrm{Shv}(G),*\bigr)$ 的点是范畴化的特征标，即特征标层：
群（叠）同态 $G\to B\G_m$。这些等价于乘性线丛 $L\to G$，即
$L_g\otimes L_h\cong L_{gh}$，或 $\mu^*L=L\boxtimes L$。所以
\[
G^\vee[1]\cong \Hom_{\mathrm{Grp}}(G,B\G_m)=\{\text{乘性线丛}\}.
\]
这完全类似于我们之前做 Cartier 对偶的方式：用 $\Spec\bigl(\C[G],*\bigr)$
定义对偶（只是函数，未做范畴化），并注意到这是群代数。

\begin{remark}
这个移位对偶发生在仿射概形范畴上阿贝尔群层构成的导出范畴中——通常的导出范畴，
但处于更万有的意义下。群概形与群叠定义这个范畴中的对象。由于该范畴是三角范畴，
在这里取移位才有意义；而移位 Cartier 对偶、即我们称为 $G^\vee[1]$ 的那个对象，
正是在这种三角意义下移位的。
\end{remark}
\srcpage{42}

还可以这样看：
\[
\Hom_{\mathrm{Grp}}(G,B\G_m)=\Ext^1(G,\G_m),
\]
所以乘性线丛 $L\to G$ 应当对应 $G$ 被 $\G_m$ 的中心扩张。这里只需取 $L$ 的总空间，
它落在阿贝尔群概形的短正合列
\[
0\longrightarrow \G_m\longrightarrow L\longrightarrow G\longrightarrow 0
\]
之中。

存在万有线丛 $L\to G\times G^\vee[1]$，因此可以把 Fourier–Mukai 变换写成积分变换：
\[
\begin{tikzcd}
& G\times G^\vee[1] \arrow[dl,"\pi_1"'] \arrow[dr,"\pi_2"] & \\
G & & G^\vee[1]
\end{tikzcd}
\]
具体地，给定 $G$ 上的层 $\cF$，其 Fourier–Mukai 对偶是
\[
\pi_{2*}\bigl(\pi_1^*\cF\otimes L\bigr).
\]

Mukai 最初考虑的是连通阿贝尔簇 $A$，例如椭圆曲线或曲线的 Jacobian。
有对偶阿贝尔簇 $A^*$ 的概念，它就是对偶环面；但 $A^*$ 与 $A^\vee[1]$ 之间有等价，
即 $A^*$ 是 $A$ 上乘性线丛构成的阿贝尔簇。\footnote{通常只用次数为零的线丛来定义，
可以证明这已经够了。}于是又得到核 $L\to A\times A^*$ 与类似于前式的核变换，
它给出导出范畴的等价。这是 Fourier–Mukai 的原型例子；Laumon 有一篇很好的例子集
[Lau91]。\footnote{人们有时用「Fourier–Mukai 变换」泛指任何积分变换，
但函数空间上像样的算子都是积分变换，导出范畴上像样的函子同样如此。
本课中我们只把「Fourier–Mukai 变换」留给那些互换卷积与乘法的变换。}
（TODO：是这篇论文吧？）

我们迟早要过渡到导出的东西，这意味着有一些背景必须先讲。

\begin{remark}
若你关心更一般的表示而非特征标，可以范畴化「$G$ 的表示就是 $G^\vee$ 上的层」这一想法。
不过在这里，范畴化的表示是「范畴的拟凝聚层」，也就是范畴化的谱分解；
但必须想清楚如何把这件事精确化——范畴的层究竟是什么？
\end{remark}

设 $C$ 是黎曼面；则 $\mathrm{Jac}(C)\simeq\mathrm{Jac}(C)^*$，而 $\Pic(C)$
（$C$ 上线丛构成的叠）等价于 $\Pic(C)\times B\G_m$，
其中后一个 $\Pic(C)$ 是我们称为 $\Pic(C)$ 的通常的非叠对象。非叠的 $\Pic(C)$
分解为 $\Z\times\mathrm{Jac}(C)$\yizhu{原文此处印作「$\Z\times\mathrm{Jac}(C)\times B\G_m$」；既已声明是非叠对象，$B\G_m$ 因子即重复，故译文按 $\Z\times\mathrm{Jac}(C)$ 排出。}{A（原文 p.42）}；对偶时 $\Z\leftrightarrow B\G_m$，
而 $\mathrm{Jac}(C)$ 自对偶，所以 $\Pic(C)$ 也自对偶，但方式很有意思：
次数与 $\G_m$-稳定化子互换了。

\begin{example}[Betti 类域论]
这与「Betti 类域论」有关，即我们讨论过的那个版本的类域论：把 $\Pic(C)$ 送到对称幺半范畴
$\bigl(\mathrm{Loc}(\Pic(C)),*\bigr)$，而 $\Spec\bigl(\mathrm{Loc}(\Pic(C)),*\bigr)$
同构于 $\mathrm{Loc}_1(C)$。也就是说，我们有 Fourier–Mukai 型的对偶
\[
\bigl(\mathrm{Loc}(\Pic(C)),*\bigr)\xrightarrow{\ \simeq\ }
\bigl(\QCoh(\mathrm{Loc}_1(C)),\otimes\bigr).
\]
左边的特征标局部系统 $\chi_L$（类似于指数函数）与右边的局部系统 $L$
（类似于 $\delta$ 分布）互换。这个变换是范畴的等价。

一种总结方式是：$\mathrm{Loc}$ 的 1-移位 Cartier 对偶就是 $\Pic$。

可以把它写成以 $\chi\to\Pic(C)\times\mathrm{Loc}_1(C)$ 为核的积分变换，
但 $\chi$ 在 $\Pic(C)$ 方向是局部系统，在 $\mathrm{Loc}_1(C)$ 方向是拟凝聚层。
之所以如此，是因为对应的两边有两种不同种类的层论。所以它有点别扭，
但理论仍然运作良好。

取一点 $x\in C$，这给出摩天楼层，从而给出映射 $\Z\to\Pic(C)$；
在这个对偶下它变成映射 $W_x:\mathrm{Loc}_1\to B\G_m$。更一般地，
Hecke 修改 $L\mapsto L(x)$ 在 $\Pic(C)$ 上按 $\Z\cdot x$ 作用，
而在对偶下这些被等同于与线丛 $W_x$ 作张量的那些算子。
可以把 $W_x$ 看成这个 Hecke 修改的特征层。
\end{example}
\srcpage{43}

\begin{example}[de Rham 类域论]
我们也可以在 de Rham 情形做这件事。这有不同的风味：对偶的两边更对称了（太好了！）
但也更复杂了（糟糕）。

给定簇 $X$，我们引入 de Rham 空间 $X_{\mathrm{dR}}$，即 $X$ 按等价关系
「$x\sim y$ 当且仅当 $x$ 与 $y$ 无穷小地接近」取的商。形式地说，
这是 $X$ 对「$X$ 在其中的对角线的形式邻域」取的商。它不是概形，
但其函子观点定义得足够好。为什么要引入这个对象？它对在代数几何中建立微分算子很有用。

对角线形式邻域上的拟凝聚层，就是 $X$ 上带有「$x$ 与 $y$ 无穷小接近时
$\cF_x\cong\cF_y$」这一数据的拟凝聚层。结果是
$\QCoh(X_{\mathrm{dR}})\simeq\mathrm{Mod}_{D_X}$，即 $X$ 上 D-模构成的范畴。
在仿射开集 $U$ 上，这些数据是一个拟凝聚层 $\cF$，连同微分算子环 $D(U)$ 在 $\cF(U)$
上的作用。

微分算子由函数与向量场生成；$\cF$ 上已有 $\cO_X$-作用，所以 D-模结构等价于指定
向量场如何作用。这可以指定为一个满足平坦性的映射 $T_X\to\End(\cF)$：
$[\xi_1,\xi_2]$ 按 $(\xi_1\cdot)$ 与 $(\xi_2\cdot)$ 的交换子作用，此外还有 Leibniz 法则
\[
\xi f=f\xi+\xi(f).
\]
所以这与 $X$ 上的平坦联络是同一回事。

这一切都讲得很快，但今天的要点是指向 de Rham 类域论。若 $G$ 是阿贝尔群，
其 de Rham 空间 $G_{\mathrm{dR}}=G/\hat G$，即对 $G$ 中单位元的形式邻域
（与 $G$ 相伴的形式群）取商。也就是说，$g_1$ 在 $G$ 中与 $g_2$ 无穷小地接近，
当且仅当 $g_1g_2^{-1}$ 在 $G$ 中与单位元无穷小地接近（或者说 $g_1g_2^{-1}\in\hat G$）。

有了这套形式，就可以回头给我们的对偶加些调料。例如对向量空间 $V$ 做通常的
Cartier 对偶，我们在对偶侧得到 $V^*$ 的形式完备化；这正说明了「代数几何里没有指数映射，
除非在形式意义下工作」这一事实。

转到 1-移位 Cartier 对偶时，$V^\vee[1]\simeq B\hat V^*$，这有点烦人：
那个形式性的问题仍然在那里。类似地，$\hat V^\vee[1]\simeq BV^*$。

现在耍个花招：对 $V$ 施 de Rham 构造，$V_{\mathrm{dR}}=V/\hat V$，
在导出范畴里它同伦等价于链复形 $\hat V\to V$。现在做移位 Cartier 对偶，
$V^\vee_{\mathrm{dR}}[1]$ 就是 $\hat V\to V$ 的移位 Cartier 对偶，于是发生这些事：
\begin{itemize}
\item $\hat V$ 映到 $V^*$；
\item $V$ 映到 $\hat V^*$；
\item 而对偶是反变的，所以我们得到映射 $V^*\leftarrow\hat V^*$。
\end{itemize}
换言之，$V_{\mathrm{dR}}$ 的 1-移位 Cartier 对偶是 $V^*_{\mathrm{dR}}$。

另一个可以看的例子是
\begin{align*}
D_{\A^1}&=\C\langle x,\partial_x\rangle\big/\bigl(\partial_x x-x\partial_x=1\bigr),\\
D_{(\A^1)^*}&=\C\langle y,\partial_y\rangle\big/\bigl(\partial_y y-y\partial_y=1\bigr).
\end{align*}
对偶把 $x\mapsto\partial_y$、$\partial_x\mapsto-y$；换言之，$D_{\A^1}$ 有一个四阶自同构，
即 Fourier 变换！它并不完全是对合，因为其平方是 $-1$。

这确实是 Fourier 变换：若 $\phi$ 是分布，它定义了一个 $D_{\A^1}$-模
\[
M_\phi:=\{L\cdot\phi\mid L\in D_{\A^1}\}.
\]
把这个变换作用到 $M_\phi$ 上，就得到 $D_{(\A^1)^*}$-模 $M_{\hat\phi}$。

Laumon–Rothstein（TODO：引文）证明了可以对 Jacobian 做这件事：若 $A=\mathrm{Jac}(C)$，
则 $A^\vee_{\mathrm{dR}}[1]$ 是 $C$ 上秩一平坦联络构成的空间。
其对偶典范地是 $\mathrm{Jac}(C)$ 上的 $\A^1$-丛。

这是最后一个阿贝尔例子；周四我们进入非阿贝尔的世界。
\end{example}
